\documentclass[11pt,a4paper]{article}
\usepackage[utf8]{inputenc}
\usepackage[T1]{fontenc}
\usepackage[margin=2.5cm]{geometry}
\usepackage{enumitem}
\usepackage{parskip}
\usepackage[hidelinks]{hyperref}

\setlist{nosep,leftmargin=1.4em}
\newcommand{\owner}[1]{\textbf{[#1]}}

\title{\vspace{-1.5em}Talking to companies: process, what we send, and who\\[-0.2em]
\large Working note for the ``spreken met bedrijven'' workstream (IOF)}
\author{Robin Wydaeghe \quad (feedback very welcome, esp.\ on contacts: \owner{Luc} \owner{Wout})}
\date{Draft, June 2026}

\begin{document}
\maketitle
\vspace{-1.5em}

\section{Purpose}
The goal of this workstream is to gauge the market, gather a few \textbf{Letters of Intent (LOIs) or
letters of support}, and get \textbf{business-plan validation} from companies, as evidence of market pull
for the IOF ``June'' call (deadline 3 August). This note covers (i) what TechTransfer requires, (ii) what we
send companies and how, (iii) draft emails, and (iv) a first list of contacts in three lanes with who owns
each intro. It is a draft written to be corrected, the main gap being the contact addresses.

\section{What TechTransfer requires (per A.\ Biondi, 26 May)}
The IP currently sits with UGent, so TechTransfer (TT) runs the confidentiality side. The agreed process:
\begin{itemize}
  \item \textbf{NDA via TT.} The mutual NDA template on the TT SharePoint is the right one. We do not fill in or send it ourselves.
  \item \textbf{Align disclosure first.} Before anything goes out, we agree with TT \emph{what} is disclosed and \emph{in which format}.
  \item \textbf{TT runs the signatures.} We send TT a \textbf{list of people to contact}. TT completes the NDA per recipient, sends it, and collects signatures.
  \item \textbf{Disclose only after signing.} We share information only once \textbf{all} relevant NDAs are signed, and only to the agreed level of detail.
  \item \textbf{Patent gate.} TT asks that \textbf{drafting of the patent application has started} before \emph{any} external sharing, in any format, NDA included. Drafting runs in parallel now and filing is targeted for $\sim$1--15 July.
\end{itemize}

\section{Timeline fit and the deadline tension}
TT estimates the full NDA cycle at \textbf{1--2 months}, which collides with the 3 August IOF deadline. The
sequence resolves this if we start now, because the early steps disclose nothing:
\begin{itemize}
  \item \textbf{Send TT the contact list this week.} Handing over a list of names is not a disclosure, so the NDA clock and the patent drafting both start immediately and run in parallel.
  \item By the time the NDAs come back signed ($\sim$July), the patent draft is under way or filed, so the gate is satisfied exactly when we need to disclose.
  \item \textbf{Keep the target set small.} Fewer NDAs sign faster, and we only need a handful of letters. Do not let the batch wait on one slow legal department.
  \item Pre-agree a single \textbf{teaser} and a single \textbf{NDA deck} with TT, so there is no back-and-forth on scope.
  \item \textbf{Backstop.} A genuinely high-level, non-confidential call needs no NDA at all. We can gauge interest and possibly secure a letter of support without ever showing the technical deck. \owner{Filip}: would such letters of support (rather than signed commercial LOIs) satisfy the IOF jury's ``market pull'' bar? The answer sets how hard we push the NDA track before August.
\end{itemize}

\section{What we send companies, and how}
\paragraph{Two documents.}
\begin{itemize}
  \item \textbf{(A) Non-confidential one-pager (teaser).} The problem, what AEGIS does (per-point absorbed power density on a 3D body in milliseconds, across the 6 to 100\,GHz mmWave-APD compliance range, with the closed-form formulation valid up to 300\,GHz), the framing of a fast pre-screen that feeds full-wave validation, two viewer screenshots, and a credentials line. No method internals. Cleared with TT before use.
  \item \textbf{(B) Confidential deck (NDA only, in meeting).} This is the two-concept deck already with Alessandro (pseudo-Brewster compensation and surface confinement, each walked through with its key figure), plus the validation figures and the coherent-MIMO operator. Shown only under a signed NDA and only to the agreed detail. Never attached to an email.
\end{itemize}

\paragraph{How we describe the method (so we stay consistent and protect the patent).}
Lead on what is genuinely and defensibly ours: a \textbf{closed-form, solver-free} evaluation of absorbed
power density over the full body surface in \textbf{milliseconds}, validated analytically (lossy-sphere Mie)
and against published numerical results to within a few percent, positioned as a \textbf{complement} that
pre-screens worst cases ahead of full-wave tools rather than replacing them. On differentiability, reframe
rather than avoid: existing differentiable-dosimetry work (Split, since 2021) wraps autodiff around a
neural-net surrogate of a numerical solver, so the gradient is approximate and slow. AEGIS's gradient is
\textbf{exact and instantaneous} because it wraps a closed-form forward model. That distinction is the
defensible part of the story. The transmission-coefficient relationship is textbook and not claimed. The
real edge is the closed-form speed plus the surface formulation, with the pseudo-Brewster compensation and
the coherent-MIMO operator staying behind the NDA.

\paragraph{Approach.} A warm introduction beats a cold email every time, so we route through the intros in
the lists below wherever they exist. Lead with credibility, keep the tone collaborative and non-commercial
(``I would value your perspective''), make one specific ask, and do not overclaim. The LOI or letter ask
comes at the \emph{end} of a call, once they have seen the demo, never in the first email. No mass emailing
and no public marketing site before the patent is filed.

\section{Draft emails}
\subsection*{Template 1: discovery / letter of support (test lab, regulator, friendly contact)}
\begin{quote}\itshape
Subject: Fast APD pre-compliance method, would value 20 minutes\\[0.3em]
Dear [name],\\[0.3em]
I am Robin Wydaeghe, a final-year PhD with Wout Joseph at UGent (INTEC-WAVES). I work on EM exposure assessment for 5G and 6G, with first-author work in npj Wireless Technology (2026) and two BioEM prizes.\\[0.3em]
Over the last months I have built a method that computes absorbed power density on a 3D body surface in milliseconds, across the 6 to 100\,GHz mmWave-APD range, validated against published numerical results to within a few percent. It is meant as a fast pre-screening layer ahead of full-wave tools, not a replacement.\\[0.3em]
I am trying to understand whether this is useful to teams like yours, and would value 20 to 30 minutes to show a short demo and hear how [APD pre-compliance / site compliance] works on your side. A one-page summary is attached.\\[0.3em]
Would a short call in the next two weeks work? Best, Robin
\end{quote}

\subsection*{Template 2: ZMT / Sim4Life (peer, collaboration, not a sales pitch)}
\begin{quote}\itshape
Subject: Closed-form APD method, possible matched validation against Sim4Life\\[0.3em]
Dear [Sven / name],\\[0.3em]
I am Robin Wydaeghe, PhD with Wout Joseph at UGent. You may know our group through the IEC TC\,106 and 63195 work. I previously placed in a Sim4Life competition with a large automation pipeline (GOLIAT), so I work inside your ecosystem rather than against it.\\[0.3em]
I have developed a closed-form surface formulation for absorbed power density at mmWave that runs in milliseconds and matches published numerical results to a few percent. I see it as a fast pre-screening complement to Sim4Life: find the worst-case beam, grip and position scenarios in seconds, then run those rigorously in Sim4Life.\\[0.3em]
Before anything else I would like to do a clean matched validation on a standard 63195 reference scenario. If the figures hold up, I think there is a natural joint contribution, possibly co-presented at BioEM. Would you be open to a short call? Best, Robin
\end{quote}

\section{Target contacts (first attempt, please correct)}
Addresses are to be sourced or confirmed by the owner. I have not guessed personal email addresses. Three
lanes, and they pull in different directions on purpose. The \textbf{fast track} is the friendly,
low-stakes end where a contact can give a letter cheaply before August. The \textbf{relationship track} is
the patient, post-patent set that should \emph{not} be rushed for an August LOI. The \textbf{academic and
institutional lane} provides adoption-signal letters that complement the industry letters for the IOF jury.
Per Wout, Sven K\"uhn is the spin-off and partnership lane while Azadeh Peyman is kept for the competition
and jury side, so those two stay in separate lanes.

\subsection*{Fast track (IOF letters before 3 August)}
\begin{itemize}[itemsep=0.7em, leftmargin=1.4em]
  \item \textbf{Ericsson Research} (Davide Colombi, Christer T\"ornevik) \quad \owner{Wout}\,\owner{Luc}\\
  Base-station and massive-MIMO actual-max EMF compliance. The strongest LOI lead in the brief. Warm path: Wout to Luc Martens to both, 2-hop.

  \item \textbf{Test labs} (CETECOM in Essen, Verkotan in Oulu, Eurofins E\&E in DE) \quad \owner{Robin}\\
  Device pre-compliance. CETECOM and Verkotan are the real targets: both contribute to IEC TC\,106 and IEEE ICES, and both will have to qualify new APD test procedures as devices above 6\,GHz ship in volume. Eurofins is the brand-name fallback. K\"uhn interfaces with all three.

  \item \textbf{BIPT (Belgium)} \quad \owner{Wout}\\
  National telecom regulator. Credentialing letter from a national authority, not a customer. Wout knows the BIPT EMF engineer. Locally easy and jury-weighty.

  \item \textbf{ARPANSA (Australia)} \quad \owner{Robin via Carolina}\\
  Australian radiation-protection authority. Worth a row only if Carolina has a real tie through her grant. Drop if she does not.

  \item \textbf{Orange Innovation} (Dinh-Thuy Phan-Huy) \quad \owner{Robin}\\
  EMF-aware beam selection and RIS research, technically tangent to our coherent-MIMO operator. The strongest operator target, both technically and politically. Warm path: via Margot Deruyck (INTEC-WAVES), 2nd-degree.

  \item \textbf{Proximus, Telenet (BE operators)} \quad \owner{Luc}\,\owner{Wout}\\
  The real wedge: Brussels-Capital Region has EMF limits much stricter than ICNIRP, with Wallonia and Flanders set differently again, so Belgian operators face a per-site compliance problem that is genuinely harder than peers. Likely letter source via Luc's standing relationship. Caveat: the actual buyer of fast site-planning tools is more often the RF engineering consultancy than the operator, so do not over-index here.

  \item \textbf{Key stakeholders named in the IDF} \quad \owner{Robin}\,\owner{Luc}\,\owner{Wout}\\
  Already identified as relevant. Extract from the IDF and confirm addresses.
\end{itemize}

\subsection*{Relationship track (post-patent, do not rush for an August letter)}
\begin{itemize}[itemsep=0.7em, leftmargin=1.4em]
  \item \textbf{ZMT / Sim4Life} (Sven K\"uhn, with N.\ Kuster, E.\ Neufeld) \quad \owner{Wout}\\
  Sim4Life ecosystem. The arc is matched validation, then co-present, then integrate. Wout endorsed Sven for the spin-off lane. Warm path: via IT'IS and IEC TC\,106.

  \item \textbf{SPEAG / IT'IS (Z\"urich)} \quad \owner{Wout}\\
  Measurement and tissue-data orbit, standards reach. Partnership conversation, not a letter.

  \item \textbf{Nokia} \quad \owner{Luc}\,\owner{Wout}\\
  RAN-vendor compliance. Warm path: via Jafar Keshvari (IEEE ICES Chair), 2nd-degree through Luc and Mike Wood. First ask is who owns EMF at Nokia.

  \item \textbf{Device side: Qualcomm, Apple, Samsung Mobile, MediaTek} \quad \owner{Wout}\\
  Post-patent only. The natural APD compliance buyer in the long run, since handset beam steering above 6\,GHz is the dominant APD use case. Route via Wout's IEEE ICES and IEC TC\,106 contacts, not cold. \textbf{Qualcomm needs disclosure-timing care}: the Hochwald exposure-aware-beam-selection patent (US11940477) sits there and is the closest industry IP to our coherent-MIMO operator. Apple is functionally unreachable cold. Samsung Mobile and MediaTek are realistic only with a named ICES contact (see request to Wout below).
\end{itemize}

\subsection*{Academic and institutional users (adoption signal, complements commercial LOIs)}
A letter from a prestigious research group saying ``we would integrate this into our antenna-EMF design loop''
is genuine market-validation evidence for the IOF jury, but it reads as \emph{adoption signal} rather than
commercial pull, so it complements rather than substitutes for the industry letters above. Two cautions for
this lane: (i) avoid the dosimetry-specific competitors (IT'IS leadership already handled, plus Hirata at
Nagoya, parts of Wiart's group at T\'el\'ecom Paris, and Poljak / Kapetanovi\'c at Split), and (ii)
disclosure discipline here is genuinely lower than in industry, so for this lane we share the teaser plus
the password-protected demo, \emph{not} the NDA deck.

\textit{Non-competing university groups (antenna, communications, RIS):}
\begin{itemize}[itemsep=0.5em, leftmargin=1.4em]
  \item \textbf{EPFL} (Anja Skrivervik, antenna and in-body) \quad \owner{Robin}\\
  High prestige, antenna-EMF design loop fit, EuCAP/AP-S community.

  \item \textbf{Lund EIT/LTH} (Mats Gustafsson, antenna bounds) \quad \owner{Robin}\\
  The provable-bounds angle. Verified email.

  \item \textbf{TUM} (Thomas Eibert, near-field) \quad \owner{Robin}\\
  Near-field rigor and validation co-author.

  \item \textbf{KTH} (Emil Bj\"ornson, massive MIMO) \quad \owner{Robin}\\
  Natural partner on the coherent-MIMO operator framing.

  \item \textbf{Paris-Saclay / CentraleSup\'elec} (Marco Di Renzo, RIS) \quad \owner{Robin}\\
  Dual-purpose with the Orange Innovation row (Phan-Huy is a co-author). Warm path: via Margot Deruyck.
\end{itemize}

\textit{Institutional and national research labs (semi-commercial, can write industry-style letters):}
\begin{itemize}[itemsep=0.5em, leftmargin=1.4em]
  \item \textbf{IETR Rennes} (Maxim Zhadobov, Giulia Sacco, mmWave APD) \quad \owner{Robin}\\
  Real APD-mmWave fit. Also a gateway to neutralising the Split prior art through co-authorship.

  \item \textbf{CEA-LETI (Grenoble)} \quad \owner{Robin}\\
  French national RF systems lab.

  \item \textbf{CNR-IEIIT (Turin)} \quad \owner{Robin}\\
  Dosimetry side of the RIS thread (already surfaced in the Gemini results).
\end{itemize}

\textit{EU consortia and testbeds (institutional weight, no competitive risk):}
\begin{itemize}[itemsep=0.5em, leftmargin=1.4em]
  \item \textbf{GOLIAT consortium} \quad \owner{Robin}\\
  We are already inside it. An internal consortium letter is essentially free and carries EU-project weight.

  \item \textbf{Oulu 6G Flagship} \quad \owner{Robin}\\
  6G testbed group, needs exposure-assessment tooling for its own experiments. Clean ``would you use this?'' pitch.
\end{itemize}

\section{What I need from you}
\begin{itemize}
  \item \owner{Wout}\,\owner{Luc}: corrections and missing addresses for the lists above. This is the main gap. In particular the BIPT engineer's name (Wout), the Nokia EMF owner via Keshvari, a green light to use the Ericsson intro through Luc, and (if it exists) a named Samsung Mobile or MediaTek contact through IEEE ICES that could move device-side onto the fast track.
  \item \owner{Wout}: availability (which days work and which do not) for the company meetings you want to join.
  \item \owner{Filip}: does the IOF jury accept letters of support and documented interest, or does it need signed commercial LOIs? And how does the jury weight academic and consortium adoption letters relative to industry LOIs? Both answers set how hard we push the NDA track before 3 August and how much we lean on the academic lane.
  \item \owner{TT (Alessandro)}: agree the disclosure scope (one cleared teaser A and one NDA-only deck B), and confirm the per-recipient NDA timing once the contact list is in.
\end{itemize}

\end{document}
